% ============================================================================
% 2. Related Work (target ~0.7 page)
% Keep third-person wording for every AELab paper (G-CAME, D-CLOSE, ODExAI)
% because review is double-blind.
% ============================================================================
\section{Related Work}
\label{sec:related}

\paragraph{Explaining object detectors.}
Explaining a detector means explaining one detection: a box, a class and a
score. Gradient-based methods adapt class activation
mapping~\cite{selvaraju2017gradcam} to this target.
ODAM~\cite{zhao2023odam} builds instance-specific maps from the gradients of
each prediction. G-CAME~\cite{nguyen2024gcame} weights the activation map
with a Gaussian centred on the predicted box.
Perturbation-based methods treat the detector as a black box.
D-RISE~\cite{petsiuk2021drise} scores random masks by the similarity between
the masked-image detections and the target detection.
D-CLOSE~\cite{truong2023dclose} applies several levels of segmentation to the
image and combines the resulting maps, and LIME~\cite{ribeiro2016lime}
fits a local surrogate model over superpixels.
Toolkits and benchmarks compare these methods across
detectors~\cite{padmanabhan2022dext,nguyen2025odexai}.
Two points matter for us. First, every such explainer takes the detection's
box as part of its target, and some, such as G-CAME, also build a box prior
into the map. Second, explanations have seldom been used downstream. The
closest exception is ODAM-NMS~\cite{zhao2023odam}, which uses the similarity
of explanations to remove duplicate boxes. We use explanations to predict
whether a detection is correct.

\paragraph{Evaluating saliency, and its pitfalls.}
The pointing game~\cite{zhang2016excitation} and its energy-based
variant~\cite{wang2020scorecam} measure how much saliency falls inside a
reference box. Insertion and deletion~\cite{petsiuk2018rise} measure how the
score changes as salient pixels are restored or removed.
Several failure modes of this kind of evaluation are known for image
classification:
\begin{itemize}
  \item Trivial baselines score well. Zhang \etal already reported a centre
  baseline for the pointing game~\cite{zhang2016excitation}. A deliberately
  uninformative ``Fake-CAM'' scores well on standard confidence-drop
  % CHECK against the full text of Poppi et al. (the abstract does not name Fake-CAM).
  metrics~\cite{poppi2021revisiting}, and dummy explanations are not reliably
  separated from real ones~\cite{stassin2023dummy}.
  \item Removal-based scores leak information through the mask
  shape~\cite{hooker2019roar,rong2022road} and depend on the occlusion
  strategy~\cite{bluecher2024occlusion}.
  \item Metrics disagree with each other and are unstable, as shown by sanity
  checks for maps~\cite{adebayo2018sanity} and for
  metrics~\cite{tomsett2020sanity}, controlled
  settings~\cite{rao2022attribution}, meta-evaluation
  toolkits~\cite{hedstrom2023quantus,hedstrom2023metaquantus} and large-scale
  benchmarks~\cite{klein2024latec}.
\end{itemize}
For detectors, Padmanabhan \etal~\cite{padmanabhan2023sanity} apply
randomization tests and find the common evaluation metrics statistically
unreliable. None of this work examines the setting where the \emph{detector's
own box} serves as pseudo-ground truth, which is how annotation-free
detection benchmarks and explainer
selection~\cite{nguyen2025odexai,cugny2022autoxai} operate.
Our audit extends the classification findings to that setting. It adds a
model-free box prior and a box-relative insertion/deletion protocol, and
reports what the metrics imply for choosing an explainer per detection.

\paragraph{Predicting detection errors.}
MetaDetect~\cite{schubert2021metadetect} introduced meta-classification
(TP \vs FP) and meta-regression (IoU) for detections. Its features are
aggregated over the box and the candidates suppressed by non-maximum
suppression. Riedlinger \etal~\cite{riedlinger2023gradient} add gradients of
the detection loss with respect to network parameters, and show that they
carry information complementary to output-based features. The protocol has
since been carried to lidar~\cite{riedlinger2025lmd}.
Other sources of uncertainty include Monte-Carlo
dropout~\cite{miller2018dropout}, ensembles~\cite{lakshminarayanan2017ensembles},
post-hoc calibration~\cite{kuzucu2024calibration,oksuz2023saod}, and
perturbation-sensitive internal features~\cite{wilson2023safe}.
Closer to deployment, run-time monitors predict detection failures from the
detector's internal representations~\cite{yatbaz2024introspection}.
Label-free evaluation estimates a detector's accuracy from the consistency
and reliability of its predictions~\cite{pcr2025}. TP/FP classifiers on
detector outputs also guide active learning~\cite{sharma2026portable}, and
DetRefiner~\cite{okazaki2026detrefiner} refines open-vocabulary detections
post hoc with a feature-fusion transformer.
Error taxonomies~\cite{hoiem2012diagnosing,bolya2020tide} separate
localization, duplicate, class-confusion and background errors. We report
results per error type because the value of an explanation should differ
between them.
Our baselines are these output-, gradient- and augmentation-based
meta-classifiers. The question is what explanation evidence adds on top of
them.

\paragraph{Explanations as signals of failure.}
Attribution statistics have been used to flag unreliable predictions in image
classification. Examples include attribution-based
confidence~\cite{jha2019abc}, adversarial-example detection from the spread
of attributions~\cite{yang2020mlloo} or from SHAP
signatures~\cite{fidel2020shap}, and out-of-distribution detection from
attribution abnormality~\cite{chen2023gaia} or learned
heatmaps~\cite{hornauer2023heatmap}. Mohseni \etal~\cite{mohseni2019predicting}
predict steering failures from saliency maps. Closest to our goal,
Spagnolo \etal~\cite{spagnolo2024xaimaps} classify lesion detections in MRI as
TP or FP from radiomic features of their saliency maps.
We study the question for object detection on natural images. We use several
explainers, faithfulness-curve shape and cross-explainer agreement. We
compare against output- and gradient-based meta-classifiers, and use placebo
maps to test whether a gain comes from the explanation itself.

\paragraph{VLM verification and risk-controlled abstention.}
Vision-language models have been used to re-score detections.
ViLD~\cite{gu2022vild} classifies region crops with CLIP embeddings. For
referring-expression comprehension, a VLM can verify a detector's proposals
one box at a time with True/False questions~\cite{vlmverify2025}. We query Qwen2.5-VL~\cite{bai2025qwen25vl} about the
box crop and, separately, about the salient pixels. The confidence score never
appears in the prompt.
For abstention we build on conformal risk control~\cite{angelopoulos2024crc},
which has been applied to object detection in safety-critical
settings~\cite{andeol2023conformal,andeol2025sequential}, and on selective
prediction~\cite{geifman2017selective,hendrycks2017baseline}. We found no prior
work that decides \emph{which} detections deserve an expensive explanation
under a compute budget. We frame this as value-of-information acquisition.
